% Source: https://artofproblemsolving.com/community/c3h2289830_bdmo_2019_junior_category_problem
% Stable identity: aops-topic-2289830-q01
Let $d=\gcd(M,N)$ and write $M=da$, $N=db$ with $\gcd(a,b)=1$. Then $\operatorname{lcm}(M,N)=dab$, and the condition becomes
\[dab=d^2\left(a^2-b^2+ab\right),\qquad\text{that is,}\qquad ab=d\,K\quad\text{with }K=a^2-b^2+ab.\]
Since $ab>0$ and $d>0$, $K$ is a positive integer, and $K$ divides $ab$.

On the other hand, $K$ is coprime to $ab$: a prime $p$ dividing $K$ and $a$ divides $K-a^2-ab=-b^2$, hence divides $b$, contradicting $\gcd(a,b)=1$; similarly a prime dividing $K$ and $b$ divides $K+b^2-ab=a^2$, hence $a$. A positive divisor of $ab$ that is coprime to $ab$ must be $1$, so $K=1$ and $d=ab$. Therefore
\[MN=d^2ab=d^2\cdot d=d^3,\]
a perfect cube. $\blacksquare$

(Such pairs exist: $a=2$, $b=3$ gives $K=4-9+6=1$, $d=6$, $M=12$, $N=18$, with $\operatorname{lcm}(12,18)=36=144-324+216$ and $MN=216=6^3$.)
